% Part: proof-theory
% Chapter: sequent-calculus
% Section: introduction

\documentclass[../../../include/open-logic-section]{subfiles}

\begin{document}

\olfileid{pt}{seq}{int}

\olsection{Introduction}

Sequent calculi are !!{proof} systems that were introduced in the
1930s by Gerhard Gentzen. They were intended less as a practical
method for !!{proving} things. Rather, Gentzen used them in order to
be able to prove certain results about the possible structure of
!!{proof}s and their property. These kinds of questions are
exactly the questions with which proof theory is concerned.

Sequent calculi operate on \emph{sequents}, as the name suggests, and
not on !!{formula}s. A sequent is a pair of sequences or multisets of
!!{formula}s. Gentzen's original systems (\Log{LK} and \Log{LJ}) used
sequences; proof theorists now prefer to use multisets. A multiset is
like a set, except that a multiset can contain an !!{element} more
than once. Just as in sets, the order of !!{element}s does not matter,
however. So the expressions $\{!A, !A, !B\}$ and $\{!A, !B, !A\}$
describe the same multiset, but that multiset is different from $\{!A,
!B\}$ and from $\{!A, !A, !B, !B\}$, since the latter contain
different numbers of $!A$ and $!B$ than the former.

Traditionally, multisets of !!{formula}s in proof theory are written
using uppercase Greek letters. So when we write $\Gamma$, $\Delta$,
$\Pi$, $\Lambda$, or $\Theta$, we intend these to be multisets of
!!{formula}s of whatever logic we are working with. Traditionally
also, instead of writing $\tuple{\Gamma, \Delta}$ proof theorists
separate the two components by a sequent symbol; we use~$\Sequent$. 

\begin{defn}[Sequent]
A \emph{sequent} is an expression of the form
\[
\Gamma \Sequent \Delta
\]
where $\Gamma$ and $\Delta$ are finite (possibly empty) multisets of
!!{formula}s. $\Gamma$ is called the \emph{antecedent}, while
$\Delta$ is the \emph{succedent}.
\end{defn}

Let $!G$~be the conjunction of the !!{formula}s in~$\Gamma$ (taking
the empty conjunction to be~$\ltrue$) and $!D$~be the disjunction of
the !!{formula}s in~$\Delta$ (taking the empty disjunction to
be~$\lfalse$). Then the sequent $\Gamma \Sequent \Delta$ is true (in
some structure~$\Struct{M}$) iff $!G \lif !D$ is true. In classical
logic, this amounts to the same as saying that either one of the
!!{formula}s in~$\Gamma$ is false, or one of the !!{formula}s
in~$\Delta$ is true.

If $\Gamma$ is a multiset of !!{formula}s, we write $\Gamma, !A$ (or
$!A, \Gamma$) for the result of adding $!A$ to~$\Gamma$. If $\Delta$
is also a multiset of !!{formula}s, then $\Gamma, \Delta$ is the
multiset union of the two sequences---if $\Gamma$ contains $!A$ with
\emph{multiplicity}~n (i.e., it contains $n$ copies of~$!A$) and
$\Delta$ with multiplicity~$m$, then $\Gamma, \Delta$ contains $!A$
with multiplicity~$n+m$. If the multiset on one side of a sequent is
empty, we usually just leave a blank. For instance, $\Sequent !A$ is
the sequent with an empty antecendent and the succedent contains $!A$
with multiplicity~$1$.

!!^a{proof} in the sequent calculus is a tree of sequents, where
the topmost (or initial) sequents are axioms in the system considered,
and if a sequent stands below one or two other sequents, it must
follow correctly by a rule of inference in the system. We will first
consider two different sequent systems for classical logic,
$\Log{G1c}$ and~$\Log{G3c}$. Their axioms and rules are given in
\cref{pt:seq:int:tab:G1c,pt:seq:int:tab:G3c}. \oliflabeldef{fol:seq::chap}{ (Gentzen's
original system \Log{LK} is described in \olref[fol][seq][]{chap}.)}{}

\subfile{rules-G1c}
\subfile{rules-G3c}

The systems are subtly different, and these differences mean that they
have different proof theoretic properties. \Log{G1c} has axioms $!A
\Sequent !A$ with no other !!{formula}s allowed. \Log{G3c} has axioms
of the form $!C, \Gamma \Sequent \Delta, !C$ with arbitrary ``side''
!!{formula}s in $\Gamma$ and~$\Delta$, but the !!{formula}~$!C$
appearing on both sides must be \emph{atomic}. \Log{G1c} has two versions
each of the $\LeftR{\land}$ and $\RightR{\lor}$ rules whereas \Log{G3c}
has just one each. And \Log{G1c} has additional ``structural rules''
called ``contraction'' (\LeftR{\Contraction}, \RightR{\Contraction})
and ``weakening'' (\LeftR{\Weakening}, \RightR{\Weakening}).

The systems \Log{G1c} and \Log{G3c} happen to be sound and complete
for classical logic. In other words, every sequent !!{provable} in
these systems is true in every !!{structure}, and every sequent that
is true in every !!{structure} has !!a{proof}. So the question
\emph{whether} a sequent has !!a{proof} can also be answered by other
methods of logic (e.g., those of model theory). And since both systems
!!{prove} exactly the sequents true in every !!{structure}, they in
particular prove the same sequents.

Proof theory, however, is not just concerned with \emph{whether}
something is provable, but \emph{how}. Proof theorists consider
questions such as ``Can a certain rule always be avoided?,'' ``Can
this rule be added to the system without resulting in new sequents
becoming provable?,'' or ``Can !!{proof}s in one system be turned into
!!{proof}s in another?'' If the answer is ``yes'' to such a question,
the proof of this fact is often proved by induction on the complexity
of !!{proof}s, or, equivalently, by induction on the structure of
!!{proof}s. This yields not just a proof of the fact (e.g., if
\Log{G1c} proves a sequent then \Log{G3c} proves the same sequent),
but a procedure (e.g., of turning !!{proof}s in~\Log{G1c} into
!!{proof}s in~\Log{G3c}).

A central result in proof theory is the \emph{cut-elimination
theorem}. It shows that we can add the cut rule
\[
\Axiom$\Gamma \fCenter \Delta, !A$
\Axiom$!A, \Pi \fCenter \Lambda$
\RightLabel{\Cut}
\BinaryInf$\Gamma, \Pi \fCenter \Delta, \Lambda$
\DisplayProof
\]
to the sequent systems without changing which sequents are
!!{provable}. Moreover, the proof provides a method that transforms
any !!{proof} using the rules of \Log{G1c} (or \Log{G3c}) together with
\Cut{} into one that only uses the rules of \Log{G1c} (or \Log{G3c}).
In other words, this procedure ``eliminates'' the \Cut{} rule from
!!{proof}s using it; hence the name.

\end{document}
